\subsection{Kunstnerisk feedback}
Til Oberstprøven vil der være en \textbf{ekstern instruktør} specifikt til at give feedback til bandet. Det foregår typisk ved at den bandeksterne ser akten fra salen og tager noter til hver nummer. Når akten er færdig, og scenepersoner får feedback på scenen, mødes bandet med den bandeksterne, samt revyens lydperson, backstage ved et af bordene. Her giver den eksterne feedback til hvert nummer.

\subsubsection{Bandets rolle: at få feedback}
Der gælder det samme for band-feedback som for generel revy-feedback: Hvis der ikke kommenteres på noget, er det som udgangspunkt godt.
Der bliver naturligvis fokuseret på de ting der kan forbedres, da formålet et at forbedre forestillingen.
Tag det derfor ikke tungt hvis der kommer mere ris end ros.
Derudover skal den eksternes kommentarer opfattes som forslag: bandet beslutter hvad bandet gør.
Således behøver man ikke overbevise den eksterne såfremt man er uenig; man kan blot vælge ikke at implementere forslaget.\\
\textbf{Alle bandmedlemmer burde medbring noget at tage noter på.} Når der gives feedback der er specifik til éen noterer man det så man kan følge op på det.\\
\textbf{En referent tager generelle noter} så bandet kan følge op på kommentarerne og afprøve ændringer.\\
\textbf{Husk at tage snacks med til bandfeedback!}

\subsubsection{Lyd}
Det er \textbf{ekstremt} gavnligt at have lydpersonen med til band-feedback. Mange problemer om balance og niveauer kan løses bedst i samarbejde mellem musikere og lydtekniker.

\subsubsection{Den eksternes rolle: at give feedback}
\paragraph{Hvad er brugbar feedback?}
Det er en vanskelig opgave at skulle give god feedback til en Oberstprøve. Man skal huske, at stort set kun har næste eftermiddag til at implementere ændringer.
Derfor burde feedback til Oberstprøven være så konkret som mulig.
Derudover skal det helst være ændringer der kan implementers med minimal øvning.\\
\textbf{God feedback til en Oberstprøve} består af et specifikt element der kan forbedres.\\
\textbf{Optimal feedback til en Oberstprøve} indeholder også et konkret forslag til forbedring.\\
\textbf{Eksempel:} Starten af omkvædet føles fladt, fordi kontrasten til broen ikke er klar nok.
For at skabe større kontrast kunne man fjerne den ene guitarrolle fra broen, og breake på 4.\\
\textbf{Eksempel:} Guitarsolen er ikke høj nok. Gør boostet større, eller gør lydtekniker opmærksom på at justere volumen on the fly.\\
Det er stadig værdifuldt at få identificeret problemer, selv hvis de ikke kommer med et konkret løsningsforslag. Vær ikke bange for at kommentere på problemer du ikke har et forslag til at løse, måske kan musikere løse det i samarbejde med teknik.

\paragraph{Konkret metode}
Jeg vil anbefale at man som ekstern får sig en aktoversigt og en stak blank papir (ca. 1 ark pr. sang). Under hver akt udfører man følgende procedure:\\
For hver sang:
\begin{itemize}
    \item Læg papiret på den brede side og noter sangens titel øverst.
    \item Træk en vertikal linie ned gennem midten af papiret.
    \item Brug venstre side til at notere kommentarer til musikalske elementer (tempo, koordination, dynamik, energi, arrangement osv.)
    \item Brug højre side til at notere kommentarer til lyd (lyddesign, balance, niveauer)
\end{itemize}
Gennemgå noterne med band og lyd efter hver akt.

\paragraph{Vibes}
Få en liste over bandmedlemmer og instrumet fra bandboss. Således kan du tiltale folk ved navn, når der skal gives feedback. Det gør alting nemmere.\\
Snak evt. med bandboss for at høre hvordan humøret er i bandet. Hvis humøret er højt og der er masser af gåpåmod kan man bare gå til den med kritik og idéer.
Hvis bandet er trætte og pressede kan man som ekstern også hjælpe ved bare at sige søde ting til dem; nogle gange har man mere brug for moral support end precise instrukser.\\

\subsubsection{Pro tips}
\begin{itemize}
    \item Bevæg dig rundt i rummet som ekstern. Lyden er ofte forskellig i de forskellige segmenter.
    \item Det kan give mening at panorere blæs lidt til højre; de er naturligt højere i venstre side.
    \item Processér bas-signalet på en eller anden måde. Et helt tørt signal peaker nemt, og begrænser bassens mulige sustain og hvor høj bassen kan være i rummet. Mindre dynamic range $\Rightarrow$ højere mulig volumen.
\end{itemize}